\section{Twisting by invertible sheaves and relative Proj}
\label{section-twisting-and-proj}

\noindent
Let $S$ be a scheme.
Let $\mathcal{A} = \bigoplus_{d \geq 0} \mathcal{A}_d$ be a
quasi-coherent graded $\mathcal{O}_S$-algebra.
Let $\mathcal{L}$ be an invertible sheaf on $S$.
In this situation we obtain another quasi-coherent graded
$\mathcal{O}_S$-algebra, namely
$$
\mathcal{B}
=
\bigoplus\nolimits_{d \geq 0}
\mathcal{A}_d \otimes_{\mathcal{O}_S} \mathcal{L}^{\otimes d}
$$
It turns out that $\mathcal{A}$ and $\mathcal{B}$ have
isomorphic relative homogeneous spectra.

\begin{lemma}
\label{lemma-twisting-and-proj}
With notation $S$, $\mathcal{A}$, $\mathcal{L}$ and $\mathcal{B}$ as
above. There is a canonical isomorphism
$$
\xymatrix{
P = \underline{\text{Proj}}_S(\mathcal{A})
\ar[rr]_g \ar[rd]_\pi & &
\underline{\text{Proj}}_S(\mathcal{B}) = P'
\ar[ld]^{\pi'} \\
& S &
}
$$
with the following properties
\begin{enumerate}
\item There are isomorphisms
$\theta_n : g^*\mathcal{O}_{P'}(n)
\to
\mathcal{O}_P(n) \otimes \pi^*\mathcal{L}^{\otimes n}$
which fit together to give an isomorphism of $\mathbf{Z}$-graded
algebras
$$
\theta :
g^*\left(
\bigoplus\nolimits_{n \in \mathbf{Z}} \mathcal{O}_{P'}(n)
\right)
\longrightarrow
\bigoplus\nolimits_{n \in \mathbf{Z}} \mathcal{O}_P(n)
\otimes \pi^*\mathcal{L}^{\otimes n}
$$
\item For every open $V \subset S$ the diagrams
$$
\xymatrix{
\mathcal{A}_n(V) \otimes \mathcal{L}^{\otimes n}(V)
\ar[r]_{multiply} \ar[d]^{\psi \otimes \pi^*}
&
\mathcal{B}_n(V) \ar[dd]^\psi \\
\Gamma(\pi^{-1}V, \mathcal{O}_P(n)) \otimes
\Gamma(\pi^{-1}V, \pi^*\mathcal{L}^{\otimes n})
\ar[d]^{multiply} \\
\Gamma(\pi^{-1}V, \mathcal{O}_P(n) \otimes \pi^*\mathcal{L}^{\otimes n})
&
\Gamma(\pi'^{-1}V, \mathcal{O}_{P'}(n)) \ar[l]_-{\theta_n}
}
$$
are commutative.
\item Add more here as necessary.
\end{enumerate}
\end{lemma}

\begin{proof}
This is the identity map when $\mathcal{L} \cong \mathcal{O}_S$.
In general choose an open covering of $S$ such that $\mathcal{L}$
is trivialized over the pieces and glue the corresponding maps.
Details omitted.
\end{proof}
















